%
% CHAPTER: Entities
%

\chapterimage{owl.pdf} % Chapter heading image

\chapter{Entities}
\label{chap:Entities}

\begin{quote}
\begin{flushright}
\emph{We are all agreed that your theory is crazy. \\
The question which divides us is whether it is crazy enough.} \\
Niels Bohr
\end{flushright}
\end{quote}
\bigskip

The theory of nescience begins with a simple question: what is the object about which knowledge is being sought? Before one can measure ignorance, compare explanations, or evaluate scientific progress, one must first identify the entities under investigation. These entities constitute the domain to which the theory will be applied.

In scientific practice, the entities studied by researchers vary enormously. They may be physical objects, such as planets, cells, or chemical compounds; phenomena, such as combustion, infection, or economic inflation; formal structures, such as groups, algorithms, or probability distributions; or theoretical constructs introduced to organize and explain observations. What unifies these cases is not a single metaphysical nature, but their role in inquiry: they are things about which researchers attempt to construct representations, collect evidence, formulate models, derive explanations, or make predictions.

For this reason, the theory of nescience does not begin with a universal collection of all possible objects. Such a collection would be mathematically problematic and philosophically unclear. Instead, every application of the theory starts by fixing a domain-specific set of entities, denoted by $\mathcal{E}$. The elements of $\mathcal{E}$ depend on the research context. In one application, $\mathcal{E}$ may be a set of mathematical structures; in another, a collection of biological organisms, diseases, physical systems, technologies, datasets, or research topics. The theory is therefore domain-relative: it does not attempt to quantify ignorance about everything at once, but about a specified collection of entities.

This chapter introduces the notion of a research entity and clarifies the formal role played by the set $\mathcal{E}$. We first explain why $\mathcal{E}$ must be fixed relative to a domain of inquiry and why universal sets are excluded. We then distinguish several types of entities that occur in scientific practice, including knowable and unknowable entities, concrete and abstract entities, observable and non-observable entities, and real and fictitious entities. This classification is not intended as a complete ontology, but as a practical taxonomy for understanding the kinds of objects that may become targets of representation and modeling.

Finally, we introduce the concept of a research area. A research area is treated formally as a subset of $\mathcal{E}$, although in practice meaningful areas are usually unified by some scientific criterion, such as a common kind, level of organization, mechanism, method, or explanatory purpose. This allows the theory of nescience to be applied not only to individual entities, but also to broader domains of inquiry. The concepts introduced in this chapter provide the foundation for the following chapters, where entities are accessed through representations and studied through descriptions.

%
% Section: Entities
%

\section{Entities}
\label{sec:descriptions_entities}

The theory of nescience studies the limits of knowledge within a given domain of investigation. Any such analysis requires a precise specification of the objects about which knowledge may be obtained. These objects will be called \emph{entities}\index{Entity}.

\begin{definition}
A \emph{research entity}\index{Research entity} is anything about which scientists attempt to construct representations, collect evidence, formulate models, or derive explanatory or predictive statements.
\end{definition}

According to this definition, a research entity may be an object, phenomenon, structure, or concept that can become the subject of systematic investigation within a given research domain.

We assume that there exists a non-empty set of \emph{entities}\index{Entities} that we seek to understand.

\begin{notation}
The set of research entities under consideration is denoted by $\mathcal{E}$.
\end{notation}

The elements of $\mathcal{E}$ are defined relative to the particular domain of knowledge in which the theory of nescience is applied. Depending on the context, this domain may correspond to a broad discipline such as mathematics, physics, engineering, or the social sciences, or it may be restricted to a specific subarea within one of these fields.

From a formal perspective, $\mathcal{E}$ should be regarded as a fixed but generally inaccessible collection: although its elements are determined by the underlying domain of inquiry, there is in general no effective (computable) procedure for accessing, deciding membership in, or enumerating its elements. 

The set $\mathcal{E}$ contains both \emph{knowable} and \emph{unknowable} entities (see Section~\ref{sec:types_of_entities} bellow for further discussion of the types of entities allowed in the theory of nesciece). 

The possibility of using universal sets is excluded; that is, the existence of a set $\xi$ containing all possible objects cannot be assumed. The problem with universal sets is that their existence is incompatible with Cantor's theorem (see Example \ref{cantor_theorem}). Cantor's theorem proves that the power set $\mathcal{P}(\xi)$, consisting of all possible subsets of $\xi$, has more elements than the original set $\xi$. This contradicts the assumption that $\xi$ includes everything. The set $\mathcal{E}$ must be explicitly fixed and domain-scoped; that is, it must denote the set of entities under consideration in a given research area, rather than an ambiguous or open-ended collection.

\begin{example}[Cantor's theorem]\index{Cantor's theorem}
\label{cantor_theorem}
Cantor's theorem proves that for any set $A$, $d(A) < d\left(\mathcal{P}(A)\right)$\footnote{Here $d(\cdot)$ denotes cardinality, see Appendix \ref{chap:Discrete_Mathematics}.}. Consider $f: A \rightarrow \mathcal{P}(A)$, a function that maps each element $x \in A$ to the set $\{x\} \in \mathcal{P}(A)$; evidently, $f$ is injective, implying $d(A) \leq d\left(\mathcal{P}(A)\right)$. To substantiate that the inequality is strict, let's assume there exists a surjective function $g: A \rightarrow \mathcal{P}(A)$ and consider the set $B = \{ x \in A : x \notin g(x) \}$. As $g$ is surjective, there must exist a $y \in A$ such that $g(y) = B$. This, however, raises a contradiction, $y \in B \Leftrightarrow y \notin g(y) = B$. Consequently, the function $g$ cannot exist, therefore $d(A) < d\left(\mathcal{P}(A)\right)$.
\end{example}

Not all conceivable sets are acceptable, as some may give rise to paradoxes. Take, for example, Russell's paradox, which proposes a set $R$ consisting of all sets that are not members of themselves. The paradox arises when we try to discern whether $R$ is a member of itself (see the Example \ref{ex:russell_paradox}). To avoid such problems, the theory of nescience is based on the Zermelo-Fraenkel set of axioms\index{Zermelo-Fraenkel set of axioms}, along with the Axiom of Choice\index{Axiom of Choice} (or ZFC). The Axiom of Separation\index{Axiom of Separation} (if $P$ is a property with parameter $p$, then for any set $x$ and parameter $p$ there exists a set $y=\{u \in x : P(u) \}$ that includes all those sets $u \in x$ that have property $P$) allows the use of this notation only to construct sets that are subsets of already existing sets. A more extensive Axiom of Comprehension\index{Axiom of Comprehension} (if $P$ is a property, then there exists a set $y=\{u : P(u) \}$) would be required to allow sets like the one proposed by Russell's paradox. Russell's paradox arises from the use of an unrestricted comprehension principle. In the axioms of ZFC, and in the theory of nescience, the axiom of comprehension is considered false.

\begin{example}[Russell's paradox]\index{Russell's paradox}
\label{ex:russell_paradox}
Suppose $R$ is the set of all sets not members of themselves, such that $R = \{ x : x \notin x \}$. The contradiction arises when querying if $R$ is a member of itself. If $R$ is not a member of itself, by its own definition, it should be; conversely, if $R$ is declared to be a member of itself, its definition dictates it should not be. Symbolically, this can be written as $R \in R \Leftrightarrow R \notin R$.
\end{example}

%
% Section: Types of Entities
%

\section{Types of Entities}
\label{sec:types_of_entities}

In scientific practice, researchers work with a wide variety of entities: physical systems, measurable phenomena, theoretical constructs, and formal objects. Some of these entities can be directly observed, others are inferred indirectly from experimental effects, and still others exist only within formal frameworks. Figure~\ref{classification-research-entities} provides a practical classification of research entities that will be useful throughout the book. This classification is not intended to be philosophically exhaustive or mathematically rigorous. Its purpose is illustrative rather than foundational. The boundaries between categories are not sharp, and in real scientific practice they often overlap or evolve over time. The diagram should therefore be read as a working taxonomy that helps clarify the kinds of objects that can become subjects of representation and modeling within the theory of nescience, rather than as a definitive ontological scheme.

\begin{figure}[t]
\centering
\begin{tikzpicture}

  % Define the ovals
  \node[draw, ellipse, minimum width=2.5cm, minimum height=1cm, align=center] (knowable) {Knowable};
  \node[draw, ellipse, minimum width=2.5cm, minimum height=1cm, right=of knowable, xshift=0.8cm] (nonknowable) {Unknowable};
  \node[draw, ellipse, minimum width=2cm, minimum height=1cm, below=of knowable, xshift=-1.2cm] (concrete) {Concrete};
  \node[draw, ellipse, minimum width=2cm, minimum height=1cm, below=of knowable, xshift=1.2cm] (abstract) {Abstract};
  \node[draw, ellipse, minimum width=1.5cm, minimum height=1cm, below=of concrete, xshift=-1.8cm] (observable) {Observable};
  \node[draw, ellipse, minimum width=1.5cm, minimum height=1cm, below=of concrete, xshift=1.8cm] (nonobservable) {Non-Observable};
  \node[draw, ellipse, minimum width=2cm, minimum height=1cm, below=of nonknowable, xshift=-1.2cm] (real) {Real};
  \node[draw, ellipse, minimum width=2cm, minimum height=1cm, below=of nonknowable, xshift=1.2cm] (ficticious) {Ficticious};

  % Arrows indicating subset relationships
  \draw[<-] (concrete) -- (knowable);
  \draw[<-] (abstract) -- (knowable);
  \draw[<-] (observable) -- (concrete);
  \draw[<-] (nonobservable) -- (concrete);
  \draw[<-] (real) -- (nonknowable);
  \draw[<-] (ficticious) -- (nonknowable);    

\end{tikzpicture}
\caption{Classification of Research Entities}
\label{classification-research-entities}
\end{figure}

% Knowable entities
\subsection*{Knowable entities}

We distinguish between knowable and unknowable entities. A \emph{knowable entity}\index{Knowable entity} is one for which, at least in principle, a sufficiently accurate scientific representation can be constructed. This does not mean that the entity is fully understood or perfectly described, but that it is accessible to systematic investigation. Scientists can refine its representations by incorporating new empirical data, correcting incorrect 
assumptions, or removing irrelevant information. Over time, this iterative process reduces error in the encoding of the entity and improves the predictive and explanatory power of the associated models.

Knowable entities can be either concrete or abstract. \emph{Concrete entities}\index{Concrete entity} are physical systems or phenomena. They exist in space and time and can interact with instruments or other physical systems. These interactions allow researchers to design experiments, collect data, and test hypotheses about their properties. Even when the entity itself is not directly accessible, its physical effects can be detected through controlled interventions and measurement protocols.

Concrete entities can be divided into observable and non-observable. \emph{Observable entities}\index{Observable entity} can be directly detected, either by the naked eye or with instruments that extend our senses, such as microscopes, telescopes, sensors, or imaging devices. Observability does not require unaided perception, but operational accessibility: the entity must produce stable signals that can be recorded, quantified, and compared across experiments. This operational criterion is what allows observable entities to serve as reliable anchors for empirical validation and model testing.

\begin{example}
A bacterial colony growing on a Petri dish is a knowable, concrete, observable entity. It can be visually inspected, counted, measured, and experimentally manipulated. Researchers can quantify its growth rate, analyze its morphology, test its response to antibiotics, or sequence its genetic material. Independent laboratories can reproduce the same experimental conditions and compare results, allowing alternative descriptions of the colony's behavior to be evaluated and refined. Because it produces stable and measurable effects (such as changes in size, color, density, or metabolic activity) its properties can be encoded in increasingly precise representations.
\end{example}

Some concrete entities are \emph{non-observable}\index{Non-observable entity}, that is, they cannot be directly observed, yet their existence can be inferred from stable, reproducible experimental effects. In such cases, the entity does not appear itself in measurement data; rather, it is postulated as the best explanation for a consistent pattern of observations. Its properties are reconstructed indirectly through mathematical models, statistical analysis, or controlled interventions.

\begin{example}
An electron is a knowable, concrete, non-observable entity. We do not observe electrons directly, but we detect their presence through cloud chamber tracks, spectral lines, scattering experiments, and electrical measurements. Their charge, mass, and spin are not accessed by direct perception, but are inferred from highly reproducible quantitative effects, such as discrete energy levels in atoms, interference patterns, and precisely measured deflections in electromagnetic fields. Competing theoretical descriptions of electrons are evaluated by comparing predicted outcomes with experimental data. Although the electron itself is never seen as a localized object in the classical sense, the stability and precision of the phenomena associated with it allow researchers to construct progressively refined representations.
\end{example}

\emph{Abstract entities}\index{Abstract entity} do not have a physical presence. They are defined within formal systems or theoretical frameworks. Although they do not interact physically with instruments, they mediate our understanding of concrete systems by offering compact and manipulable representations. Equations, probability distributions, algorithms, state spaces, and optimization functions allow researchers to formalize assumptions, derive consequences, and compute quantitative forecasts that can later be confronted with empirical data. In many cases, progress in a discipline results from introducing a new abstract entity that reorganizes existing observations under a more coherent structure, reduces descriptive complexity, or increases predictive accuracy. 

\begin{example}
The notion of a Turing machine (see Definition \ref{def:Turing-Machine}) is a knowable, abstract entity. It is not a physical object, but it is precisely defined, can be formally analyzed, and is used to prove theorems about computability and complexity. Its components are specified with exact rules, allowing rigorous reasoning about what can and cannot be computed. Although no physical device perfectly realizes the infinite tape assumed in the definition, the abstraction captures essential structural properties shared by all digital computers. Because its behavior is fully determined by formal rules, competing representations can be unambiguously compared, and proofs about decidability, universality, and incompleteness can be carried out with mathematical precision.
\end{example}

% Unknowable entities
\subsection*{Unknowable entities}

An \emph{unknowable entity}\index{Unknowable entity} is an entity for which no reliable and stable representation can be constructed within any scientific framework. In this sense, unknowability is structural rather than temporary: it does not arise from current technological limitations or lack of data, but from the absence of a possible operational connection between the entity and reproducible procedures. An entity is structurally unknowable when no experiment, measurement protocol, or formal derivation could, even in principle, discriminate between alternative representations or progressively reduce error in its encoding. Under this definition, unknowable entities fall outside the domain of scientific research.

Unknowable entities can be classified as real or fictitious. Some entities are \emph{real}\index{Real entity} in the sense that they belong to physical reality, yet remain structurally beyond our capacity to investigate or represent with scientific precision. Even if the entity exists independently of our knowledge, no reproducible procedure could in principle establish reliable constraints on its properties. As a consequence, scientific representations cannot converge, and competing encodings cannot be meaningfully evaluated. These entities may be logically conceivable and even physically instantiated, but they remain outside the domain of entities that can be stabilized through empirical or formal refinement.

\begin{example}
Consider the exact microphysical state of the universe at a precise instant of time, including the exact position and momentum of every particle with infinite precision. If physical reality is continuous (as assumed in standard classical models), then such a state is a real entity: it corresponds to an actual configuration of the world. However, it is structurally unknowable. No finite physical system can measure, store, or encode an infinite amount of information with arbitrary precision. Moreover, no experimental protocol could discriminate between two candidate representations that differ only beyond any physically realizable resolution. Therefore, although the entity may exist, no reliable and stable scientific representation of it can in principle be constructed. 
\end{example}

\emph{Ficticious entities}\index{Ficticious entity} are introduced hypothetically but do not correspond to any reproducible phenomenon. Such entities may arise as speculative explanatory devices, heuristic constructions, or artifacts of incomplete theories. Unlike real but unknowable entities, fictitious entities do not merely escape representation due to structural limits; rather, they fail to anchor representations to observable effects or formal necessities. They are usually abandoned when alternative frameworks explain the same phenomena without invoking them, or when continued experimental effort fails to produce reproducible consequences associated with their existence.

\begin{example}
The luminiferous ether, once postulated as the medium for light propagation, is an example of an unknowable and fictitious entity. It was introduced in the nineteenth century to account for the apparent need of a mechanical medium through which electromagnetic waves could travel, in analogy with sound waves in air. However, despite increasingly precise experimental efforts to detect the Earth's motion relative to this supposed medium, no consistent or reproducible effect was ever observed. Eventually, the development of special relativity rendered the ether unnecessary by providing a coherent theoretical framework that explained electromagnetic phenomena without invoking any background medium.
\end{example}

%
% Research Areas
%

\section{Research Areas}
\label{sec:research_areas}

Scientific inquiry is rarely confined to a single entity. Researchers usually study collections of entities that belong to a common domain, share relevant properties, appear at a similar level of organization, or participate in a common explanatory framework. Such collections will be called \emph{research areas}. The concept is useful because many scientific questions are not directed at isolated entities, but at families of entities considered together: biological species, chemical compounds, mathematical structures, diseases, algorithms, economic systems, or technological artifacts.

\begin{definition}
Let $\mathcal{E}$ be a set of entities. We define a \emph{research area}\index{Research area} as a subset of entities
\[
\mathcal{A} \subseteq \mathcal{E}.
\]
\end{definition}

The definition is intentionally simple. A research area is formally treated as a subset of the set of entities under consideration. However, in scientific practice, not every subset of $\mathcal{E}$ constitutes a meaningful research area. Scientifically useful areas are usually unified by some intelligible criterion. The entities in an area may belong to the same kind (see Section XX), appear at the same level of organization, participate in the same mechanism, share relevant attributes, be studied by similar methods, or contribute to the same explanatory objective. The exact criterion depends on the domain in which the theory of nescience is applied.

\begin{example}
If $\mathcal{E}$ is the set of entities studied in biology, then possible research areas include mammals, viruses, metabolic pathways, cancer types, or ecosystems. These areas are not arbitrary collections. Mammals are grouped by biological classification; viruses by their biological organization and replication mechanisms; metabolic pathways by their functional role in cellular processes; cancer types by pathological and molecular criteria; and ecosystems by the interaction of organisms with their environment. Each case reflects a different principle of scientific grouping.
\end{example}


In Chapter~\ref{chap:Descriptions}, we will study how a research areas can be characterized through descriptions that jointly account for the representations of its members. From this perspective, a research area is not merely a set of entities selected in advance, but may also be understood through the common descriptive structure that makes those entities scientifically intelligible as a group.

Research areas may also have subareas.

\begin{definition}
Let $\mathcal{E}$ be a set of entities, and $\mathcal{A} \subseteq \mathcal{E}$ a research area. We define a \emph{research subarea}\index{Research subarea} of the area $\mathcal{A}$ as a subset of entities
\[
\mathcal{B} \subseteq \mathcal{A}.
\]
\end{definition}

This definition gives rise to hierarchies of areas, in which broad domains are progressively refined into more specialized subdomains.

\begin{example}
Within the area of mammals, one may distinguish primates, rodents, carnivores, or marine mammals. Within the area of machine learning, one may distinguish supervised learning, unsupervised learning, reinforcement learning, or representation learning. 
\end{example}

It is important, however, that such hierarchies need not be unique. The same set of entities may be organized in different ways depending on the scientific purpose. A biological domain may be organized taxonomically, grouping organisms into species, genera, families, and higher biological categories. The same organisms may also be grouped ecologically, according to habitat, trophic role, or environmental adaptation. Both organizations may be scientifically legitimate, because they answer different questions.

This flexibility is important for the theory of nescience. The value of a research area depends on the purpose of the investigation. A collection of entities may be appropriate for one question and inappropriate for another. If the goal is to study evolutionary ancestry, a taxonomic area may be useful. If the goal is to study disease transmission, an area based on host-pathogen interaction may be more appropriate. If the goal is to compare predictive models, an area may be defined by the availability of common datasets, common variables, or common modeling assumptions.

Research areas therefore play a practical role. They allow the theory of nescience to be applied not only to individual entities, but also to broader domains of inquiry. Once a research area $\mathcal{A}$ has been fixed, one may ask how well the entities in $\mathcal{A}$ are represented, how accurately they are described, how much unnecessary complexity is present in their descriptions, and how much nescience remains in the area as a whole. In later chapters, these questions will lead to area-level versions of miscoding, mismodel, and nescience.

%
% Section: References
%

\section{References}

For more information about Russell's paradox, Cantor's theorem, and universal sets, refer, for example, to \cite{jech2013set}. This book also covers the Zermelo-Fraenkel system of axioms, including the Axiom of Choice. 


